% Part: first-order-logic
% Chapter: tableaux

\documentclass[../../../include/open-logic-chapter]{subfiles}

\begin{document}

\iftag{FOL}
      {\olchapter{fol}{tab}{టాబ్లోలు}}
      {\olchapter{pl}{tab}{టాబ్లోలు}}

\begin{editorial}
ఈ అధ్యాయం చిహ్నిత విశ్లేషణాత్మక టాబ్లో వ్యవస్థను పరిచయం చేస్తుంది.

నిరూపణ వ్యవస్థగా టాబ్లోలకు సంబంధించిన విషయాన్ని చేర్చడానికి లేదా
తొలగించడానికి ``prfTab'' ట్యాగును ఉపయోగించండి.
\sourcecorrection{OLTETAB-001}{ఈ టాబ్లో అధ్యాయం, దాని
\texttt{prfTab} ట్యాగు గురించి చెప్పే చోట ఆంగ్ల మూలం పొరపాటుగా
“natural deduction” అని పేర్కొంది. అధ్యాయ మార్గం, గుర్తింపులు,
దిగుమతులన్నిటికీ అనుగుణంగా దాన్ని “టాబ్లోలు”గా సరిచేశాం.}
\end{editorial}

\olimport{rules-and-proofs}

\olimport{propositional-rules}

\iftag{FOL}{%
  \olimport{quantifier-rules}
}{}

\olimport{derivations}

\olimport{proving-things}

\iftag{FOL}{%
  \olimport{proving-things-quant}
}{}

\olimport{proof-theoretic-notions}

\olimport{provability-consistency}

\olimport{provability-propositional}

\iftag{FOL}{%
  \olimport{provability-quantifiers}
}{}

\olimport{soundness}

\iftag{FOL}{%
  \olimport{identity}
  \olimport{soundness-identity}
}{}

\OLEndChapterHook

\end{document}
